\section{Innovation I: New Product Development (NPD)}
\textbf{The Success Gap:} Less than 3 out of 10 new products survive their first year.
Retail store data shows only 24\% of new items sustain sales by Week 52.

\subsection{Traditional vs. Emerging Approaches}
\begin{itemize}
    \item \textbf{Traditional Approach:} An aggregate or firm-specific linear model (R\&D
          $\rightarrow$ Production $\rightarrow$ Marketing). It is expert-driven,
          sequential, and rigid, making it slow and costly.
    \item \textbf{Emerging Approaches:} Includes \textbf{Design Thinking}, \textbf{Lean Startup}, and \textbf{Jobs-to-be-Done}.
    \item \textbf{Core Principles:} These are \textbf{Customer Centric} (start from problems/needs), \textbf{Experimental} (build experiments), and \textbf{Iterative} (feedback loops).
\end{itemize}

\subsection{Minimum Viable Product (MVP)}

An MVP is a prototype designed to test fundamental business hypotheses[cite:
        1144, 1145]. Its goal is to \textbf{begin the process of learning}, not to reach product perfection.
\begin{itemize}
    \item \textbf{Key Components:} It must have enough basic functionality to validate the
          concept and solve a specific problem.
    \item \textbf{The Innovation Paradox:} Innovators are advised to avoid large, irreversible
          investments before validating demand, but demand is often only observable after
          such investments are made.
\end{itemize}

\subsection{MVP Examples \& Prototypes}
\begin{itemize}
    \item \textbf{Dropbox:} Used a \textbf{Video MVP} and a landing page to illustrate how the product would work for
          early adopters. This validated the hypothesis that users appreciated file
          synchronization before building the technology.
    \item \textbf{Gojek:} The first prototype was a call center to match riders and customers. It
          appeared automated but was manually run, allowing them to validate demand
          before investing in automation.
    \item \textbf{Landing Pages:} Tools like Wix or Strikingly are used to communicate a \textbf{Value Proposition} and collect feedback via visits or email subscriptions.
\end{itemize}

\subsection{Strategic Experimentation}
\textbf{Falsifiable Hypothesis:} Identify the riskiest assumption using the format: "We
believe [number of users] will [behavior] within [duration]".
\begin{itemize}
    \item \textbf{Key Metrics:} Use \textbf{Unique Visitors}, \textbf{Subscribers}, and \textbf{Conversion Rate} ($\frac{\text{Subscribers}}{\text{Visitors}}$) to measure
          results.
    \item \textbf{Validated Learning:} Interpret data to confirm or reject hypotheses and propose
          the next action plan.
\end{itemize}

\textbf{Case Study: GETWEAR Failure:} Failed due to a lack of proper demand validation
before big investments and ignoring negative feedback.